% Part: propositional-logic
% Chapter: propositional-logic
% Section: preliminaries

\documentclass[../../../include/open-logic-section]{subfiles}

\begin{document}

\olfileid{pl}{syn}{pre}

\olsection{Voorbereidingen}

\begin{thm}[\emph{Inductieprincipe voor !!{formula}s}]
  \ollabel{thm:induction}  
  Stel dat een eigenschap~$P$ voor alle atomaire !!{formula}s geldt en
  bovendien aan de volgende voorwaarden voldoet:
  \begin{enumerate}
    \tagitem{prvNot}{zij geldt voor $\lnot !A$ wanneer zij voor~$!A$
      geldt;}{}
    \tagitem{prvAnd}{zij geldt voor $(!A \land !B)$ wanneer zij voor
      $!A$ en~$!B$ geldt;}{}
    \tagitem{prvOr}{zij geldt voor $(!A \lor !B)$ wanneer zij voor
      $!A$ en~$!B$ geldt;}{}
    \tagitem{prvIf}{zij geldt voor $(!A \lif !B)$ wanneer zij voor
      $!A$ en~$!B$ geldt;}{}
    \tagitem{prvIff}{zij geldt voor $(!A \liff !B)$ wanneer zij voor
      $!A$ en~$!B$ geldt.}{}
  \end{enumerate}
  Dan geldt $P$ voor alle !!{formula}s.
\end{thm}

\begin{proof}
  Zij $S$ de collectie van alle !!{formula}s met eigenschap~$P$.
  Duidelijk is $S \subseteq \Frm[L_0]$. De collectie~$S$ voldoet aan
  alle voorwaarden van \olref[fml]{defn:formulas}: zij bevat alle
  atomaire !!{formula}s en is gesloten onder de !!{operator}s.
  $\Frm[L_0]$ is de kleinste dergelijke klasse, dus
  $\Frm[L_0] \subseteq S$. Bijgevolg is $\Frm[L_0] = S$ en heeft elke
  formule eigenschap~$P$.
\end{proof}

\begin{prop}\ollabel{prop:balanced}
   Elke !!{formula} in~$\Frm[L_0]$ is \emph{in balans}: zij bevat
   evenveel linkerhaakjes als rechterhaakjes.
\end{prop}

\begin{prob}
Bewijs \olref[pl][syn][pre]{prop:balanced}
\end{prob}

\begin{prop} \ollabel{prop:noinit}
Geen echt beginstuk van !!a{formula} is !!a{formula}.
\end{prop}

\begin{prob}
  Bewijs \olref[pl][syn][pre]{prop:noinit}
\end{prob}

\begin{prop}[Unieke leesbaarheid]
Elke !!{formula}~$!A$ in $ \Frm[L_0]$ heeft precies één ontleding in
een van de volgende vormen:
\begin{enumerate}
\tagitem{prvFalse}{$\lfalse$.}{}

\tagitem{prvTrue}{$\ltrue$.}{}

\item $\Obj p_n$ voor een zekere $\Obj p_n \in  \PVar$.
  
\tagitem{prvNot}{$\lnot !B$ voor een zekere !!{formula}~$!B$.}{}

\tagitem{prvAnd}{$(!B \land !C)$ voor zekere !!{formula}s $!B$ en~$!C$.}{}

\tagitem{prvOr}{$(!B \lor !C)$ voor zekere !!{formula}s $!B$ en~$!C$.}{}

\tagitem{prvIf}{$(!B \lif !C)$ voor zekere !!{formula}s $!B$ en~$!C$.}{}

\tagitem{prvIff}{$(!B \liff !C)$ voor zekere !!{formula}s $!B$ en~$!C$.}{}
\end{enumerate}
Bovendien is deze ontleding \emph{uniek}.
\end{prop}

\begin{proof}
Door inductie naar $!A$. Stel bijvoorbeeld dat $!A$ twee verschillende
ontledingen heeft, als $(!B \lif !C)$ en als $(!B' \lif !C')$. Dan
moeten $!B$ en $!B'$ gelijk zijn; anders zou een van beide een echt
beginstuk van de andere zijn. Als de twee ontledingen van $!A$
verschillen, moet dat dus komen doordat $!C$ en $!C'$ verschillende
ontledingen van dezelfde tekenreeks zijn. Dit is volgens de
inductiehypothese onmogelijk.
\end{proof}


\begin{defn}[Uniforme substitutie]
Als $!A$ en $!B$ !!{formula}s zijn en $\Obj p_i$ !!a{propositional
  variable} is, dan duidt $\Subst{!A}{!B}{\Obj p_i}$ het resultaat aan
  van het vervangen van elk voorkomen van $\Obj p_i$ door een voorkomen
  van $!B$ in $!A$. Evenzo wordt de gelijktijdige substitutie van
$\Obj p_1$, \dots,~$\Obj p_n$ door de !!{formula}s $!B_1$,
\dots,~$!B_n$ aangeduid met
$\SSubst{!A}{\subst{!B_1}{\Obj p_1},\dots,\subst{!B_n}{\Obj p_n}}$.
\end{defn}

\begin{prob} Bepaal voor elk van de vijf onderstaande !!{formula}s of
 die !!{formula} kan worden uitgedrukt als een substitutie
 \( \Subst{!A}{!B}{\Obj p_i} \), waarbij \( !A \) gelijk is aan (i)
 \( \Obj p_0 \); (ii) \( ( \lnot \Obj p_0 \land \Obj p_1) \); en (iii)
 \( ( ( \lnot \Obj p_0 \lif \Obj p_1 ) \land \Obj p_2 ) \). Geef in
  elk geval de bijbehorende substitutie.
  \begin{enumerate}
    \item \( \Obj p_1 \) 
    \item \( ( \lnot \Obj p_0 \land \Obj p_0 ) \)
    \item \( ( ( \Obj p_0 \lor \Obj p_1 ) \land \Obj p_2 )  \)
    \item \( \lnot ( ( \Obj p_0 \lif \Obj p_1 ) \land \Obj p_2 ) \)
    \item \( (( \lnot ( \Obj p_0 \lif \Obj p_1 ) \lif ( \Obj p_0 \lor \Obj p_1 )) \land \lnot ( \Obj p_0 \land \Obj p_1 )) \)
  \end{enumerate}
\end{prob}

\begin{prob}
Geef met inductie een wiskundig precieze definitie van
$\Subst{!A}{!B}{p}$.
\end{prob}
\end{document}
